\section{Pendahuluan}
Informasi yang diperlukan pengguna sering tersedia dalam dokumen, tetapi tersebar di beberapa bagian dengan istilah dan struktur yang berbeda. RAG mengambil informasi dari sumber eksternal sebelum model bahasa membentuk jawaban \cite{lewisRetrievalAugmentedGenerationKnowledgeIntensive2021}. Keberhasilan tahap ini bergantung pada bagaimana dokumen diubah menjadi unit informasi yang dapat ditemukan kembali.

Peraturan perundang-undangan dan \tech{User Manual} memperlihatkan dua kebutuhan yang berbeda. Pasal merupakan unit struktural utama dalam representasi dokumen hukum. Pemahaman pada tingkat pasal memerlukan rincian di bawahnya, yaitu ayat, huruf, dan angka. Sementara itu, petunjuk produk tersusun atas bagian, subbagian, dan langkah prosedur. Pemotongan berdasarkan panjang teks saja dapat memisahkan rincian dari konteks yang menjelaskannya.

Strategi pencarian juga perlu disesuaikan. \tech{Sparse retrieval} memanfaatkan kecocokan istilah, sedangkan \tech{dense retrieval} menggunakan representasi vektor untuk menemukan kedekatan makna \cite{robertsonProbabilisticRelevanceFramework2009,karpukhinDensePassageRetrieval2020}. Performa pada satu koleksi belum tentu berlaku pada koleksi lain, sebagaimana menjadi perhatian dalam evaluasi lintas \tech{domain} \cite{thakurBEIRHeterogeneousBenchmark2021}. Relasi antarperaturan menambah kebutuhan untuk mempertahankan hubungan yang tidak cukup direpresentasikan sebagai potongan teks terpisah.

Penelitian ini mengembangkan \tech{pipeline} konstruksi basis pengetahuan berbasis \tech{plugin} pada Omni-RAG. Kontribusinya mencakup pembentukan \tech{chunk} dengan identitas sumber, penyediaan strategi \tech{sparse}, \tech{dense}, dan \tech{hybrid}, serta pembentukan \tech{knowledge graph} dari \tech{chunk} yang sama. Ketiganya memakai kontrak yang memungkinkan penggantian strategi tanpa mengubah orkestrator inti. Evaluasi membandingkan konfigurasi pada dua korpus untuk menilai kemampuan menemukan kembali sumber yang relevan.

\tech{Context PII}, \tech{guardrail}, dan \tech{history compactor} dibahas sebagai komponen pendukung percakapan. Kontribusi komponen tersebut tidak diukur dalam eksperimen. Akuisisi dokumen, \tech{parser}, dan pembentukan jawaban menjadi konteks integrasi sistem, sedangkan fokus artikel ini adalah konstruksi basis pengetahuan dan hasil \tech{retrieval}.
